\documentclass[aspectratio=169,10pt]{beamer}

% Shared Beamer setup for Fall 2026 course slides.
\mode<presentation>{
  \usetheme{Madrid}
  \usecolortheme{default}
}

\definecolor{CalPolyGreen}{HTML}{154734}
\definecolor{CalPolyGold}{HTML}{BD8B13}
\definecolor{DeckInk}{HTML}{1F2933}
\definecolor{DeckMuted}{HTML}{5F6B7A}

\setbeamercolor{structure}{fg=CalPolyGreen}
\setbeamercolor{palette primary}{bg=CalPolyGreen,fg=white}
\setbeamercolor{palette secondary}{bg=DeckInk,fg=white}
\setbeamercolor{palette tertiary}{bg=CalPolyGold,fg=DeckInk}
\setbeamercolor{frametitle}{bg=CalPolyGreen,fg=white}
\setbeamercolor{title}{fg=CalPolyGreen}
\setbeamercolor{normal text}{fg=DeckInk,bg=white}
\setbeamercolor{block title}{bg=CalPolyGreen,fg=white}
\setbeamercolor{block body}{bg=black!3,fg=DeckInk}

\setbeamertemplate{navigation symbols}{}
\setbeamertemplate{footline}[frame number]
\setbeamertemplate{itemize item}{\color{CalPolyGold}$\blacktriangleright$}
\setbeamertemplate{itemize subitem}{\color{CalPolyGreen}$\bullet$}

\newcommand{\CourseNumber}{Course}
\newcommand{\CourseTitle}{Course Title}
\newcommand{\LectureNumber}{0}
\newcommand{\LectureTitle}{Lecture Title}
\newcommand{\LectureDate}{Date}
\newcommand{\CourseUnit}{Unit}

\newcommand{\coursenumber}[1]{\renewcommand{\CourseNumber}{#1}}
\newcommand{\coursetitle}[1]{\renewcommand{\CourseTitle}{#1}}
\newcommand{\lecturenumber}[1]{\renewcommand{\LectureNumber}{#1}}
\newcommand{\lecturetitle}[1]{\renewcommand{\LectureTitle}{#1}}
\newcommand{\lecturedate}[1]{\renewcommand{\LectureDate}{#1}}
\newcommand{\courseunit}[1]{\renewcommand{\CourseUnit}{#1}}

\author{Kirk Duran}
\institute{California Polytechnic State University, San Luis Obispo}
\date{\LectureDate}

\newcommand{\makedecktitle}{
  \begin{frame}[plain]
    \vspace{0.25in}
    {\usebeamerfont{title}\color{CalPolyGreen}\LectureTitle\par}
    \vspace{0.18in}
    {\large \CourseNumber{}: \CourseTitle\par}
    \vspace{0.08in}
    {\large Lecture \LectureNumber{} -- \LectureDate\par}
    \vfill
    {\small Kirk Duran \textbar{} Cal Poly, San Luis Obispo\par}
  \end{frame}
}

\newcommand{\placeholdernote}{
  \begin{block}{Placeholder}
    Replace these bullets with the final lecture material, examples, and in-class activities.
  \end{block}
}

\AtBeginSection[]{
  \begin{frame}{Roadmap}
    \tableofcontents[currentsection]
  \end{frame}
}


\pdfstringdefDisableCommands{\def\translate#1{#1}}

\graphicspath{{../assets/lecture-10/}}

% If an image file is not yet available, the slide displays a labeled
% placeholder using the intended filename. Place the matching file in
% ../assets/lecture-10/ to replace the placeholder automatically.
\newcommand{\lectureimage}[2][1.75in]{%
  \IfFileExists{../assets/lecture-10/#2}{%
    \includegraphics[width=\linewidth,height=#1,keepaspectratio]{#2}%
  }{%
    \fbox{%
      \parbox[c][#1][c]{0.94\linewidth}{%
        \centering
        \small\textbf{Image Placeholder}\\[0.08in]
        \scriptsize\texttt{\detokenize{#2}}
      }%
    }%
  }%
}

\newcommand{\lectureplaceholder}[2][1.75in]{%
  \fbox{%
    \parbox[c][#1][c]{0.94\linewidth}{%
      \centering
      \small\textbf{Image Placeholder}\\[0.08in]
      \scriptsize #2
    }%
  }%
}

\setbeamersize{text margin left=0.42in,text margin right=0.42in}

\coursenumber{CS 1001}
\coursetitle{Introduction to Computing}
\lecturenumber{10}
\lecturetitle{Functions III --- Composition}
\lecturedate{Fall 2026}
\courseunit{1. Computing and Program Execution}

\begin{document}

\makedecktitle

\begin{frame}{Today}
  \begin{itemize}
    \item Extend the one-call model to programs with multiple active function calls.
    \item Represent active calls using separate execution environments organized by a call stack.
    \item Distinguish runtime execution environments from lexical scope and name visibility.
    \item Trace functions that call other functions, including nested call expressions.
    \item Use returned values to compose smaller computations into larger ones.
    \item Treat built-in and library functions through the same argument--return-value interface.
  \end{itemize}

  \begin{keyidea}
    Lecture 9 traced one active call. Today we keep unfinished callers active while additional calls begin.
  \end{keyidea}
\end{frame}

\sectionframe{Part 1: From One Active Call to Multiple Calls}

\begin{frame}{One Function Can Call Another Function}
  \begin{columns}[T,onlytextwidth]
    \begin{column}{0.54\textwidth}
      \begin{block}{Program}
        \texttt{def double(n: int) -> int:}\\
        \hspace*{0.20in}\texttt{return 2 * n}\\[0.08in]
        \texttt{def square\_perimeter(length: int) -> int:}\\
        \hspace*{0.20in}\texttt{two\_sides = double(length)}\\
        \hspace*{0.20in}\texttt{return two\_sides + two\_sides}
      \end{block}

      \begin{itemize}
        \item Calling \texttt{square\_perimeter} begins one function call.
        \item Its body then begins a second call to \texttt{double}.
      \end{itemize}
    \end{column}

    \begin{column}{0.40\textwidth}
      % Intended visual: square_perimeter calling double while the outer call remains active.
      \lectureimage[2.75in]{slide4.png}
    \end{column}
  \end{columns}
\end{frame}

\begin{frame}{The Caller Is Suspended, Not Finished}
  \begin{block}{Execution reaches}
    \centering
    \Large\texttt{two\_sides = double(length)}
  \end{block}

  \begin{enumerate}
    \item Evaluate \texttt{length} to obtain the argument value.
    \item Begin a new call to \texttt{double}.
    \item Suspend execution of \texttt{square\_perimeter} at the call site.
    \item Execute \texttt{double} until it returns.
    \item Resume \texttt{square\_perimeter} and complete the assignment.
  \end{enumerate}

  \begin{keyidea}
    An unfinished caller retains its local bindings while the callee executes.
  \end{keyidea}
\end{frame}

\begin{frame}{Active Calls Form a Stack}
  \begin{cpdefinition}
    The \textbf{call stack} is the runtime ordering of currently active function calls, organized by caller--callee relationships.
  \end{cpdefinition}

  \begin{columns}[T,onlytextwidth]
    \begin{column}{0.48\textwidth}
      \begin{block}{While \texttt{double} executes}
        \centering
        \texttt{double}\\
        $\uparrow$ current call\\[0.08in]
        \texttt{square\_perimeter}\\[0.08in]
        initial code
      \end{block}
    \end{column}

    \begin{column}{0.46\textwidth}
      \begin{itemize}
        \item The most recently started unfinished call is current.
        \item Earlier callers remain beneath it.
        \item Returning removes the current call and resumes its immediate caller.
      \end{itemize}
    \end{column}
  \end{columns}
\end{frame}

\begin{frame}{Each Active Call Has Its Own Execution Environment}
  \begin{columns}[T,onlytextwidth]
    \begin{column}{0.46\textwidth}
      \begin{block}{\texttt{double} call}
        \texttt{n -> 2}
      \end{block}

      \begin{block}{\texttt{square\_perimeter} call}
        \texttt{length -> 2}\\
        \texttt{two\_sides}: not yet bound
      \end{block}
    \end{column}

    \begin{column}{0.46\textwidth}
      \begin{itemize}
        \item Each call has a distinct set of parameter and local-variable bindings.
        \item The outer call's bindings do not disappear when the inner call begins.
        \item The current frame changes as calls begin and return.
      \end{itemize}
    \end{column}
  \end{columns}

  \begin{keyidea}
    Same program, multiple simultaneously active environments.
  \end{keyidea}
\end{frame}

\begin{frame}{Returning Unwinds One Call at a Time}
  \begin{columns}[T,onlytextwidth]
    \begin{column}{0.51\textwidth}
      \begin{block}{Inside \texttt{double}}
        \texttt{return 2 * n}\\[0.05in]
        \texttt{2 * 2 -> 4}\\
        return \texttt{4}
      \end{block}

      \begin{block}{Resume \texttt{square\_perimeter}}
        \texttt{two\_sides = 4}\\
        \texttt{return 4 + 4}\\
        return \texttt{8}
      \end{block}
    \end{column}

    \begin{column}{0.43\textwidth}
      % Intended visual: stack shrinking from double -> square_perimeter -> initial code.
      \lectureimage[2.65in]{slide8.png}
    \end{column}
  \end{columns}
\end{frame}

\sectionframe{Part 2: Execution Environments and Scope}

\begin{frame}{Execution Environment and Scope Are Related but Different}
  \begin{columns}[T,onlytextwidth]
    \begin{column}{0.46\textwidth}
      \begin{cpdefinition}
        An \textbf{execution environment} is a runtime mapping from names to values for one active call.
      \end{cpdefinition}

      \begin{block}{Runtime question}
        What bindings exist in this active call right now?
      \end{block}
    \end{column}

    \begin{column}{0.46\textwidth}
      \begin{cpdefinition}
        A \textbf{scope} determines where a name's binding is visible in program text.
      \end{cpdefinition}

      \begin{block}{Name-resolution question}
        Which binding does this occurrence of a name refer to?
      \end{block}
    \end{column}
  \end{columns}
\end{frame}

\begin{frame}[shrink=3]{The Same Spelling Can Refer to Different Bindings}
  \begin{block}{Program fragment}
    \texttt{x = 10}\\[0.06in]
    \texttt{def compute(value):}\\
    \hspace*{0.20in}\texttt{x = value + 1}\\
    \hspace*{0.20in}\texttt{return x}
  \end{block}

  \begin{columns}[T,onlytextwidth]
    \begin{column}{0.46\textwidth}
      \begin{block}{Module-level binding}
        \texttt{x -> 10}
      \end{block}
    \end{column}

    \begin{column}{0.46\textwidth}
      \begin{block}{During \texttt{compute(4)}}
        \texttt{value -> 4}\\
        local \texttt{x -> 5}
      \end{block}
    \end{column}
  \end{columns}

  \begin{keyidea}
    A variable name is not globally unique; its meaning depends on its scope and active environment. A local assignment to \texttt{x} does not rebind a separate module-level \texttt{x}.
  \end{keyidea}
\end{frame}

\begin{frame}{A Callee Does Not Inherit Its Caller's Local Variables}
  \begin{itemize}
    \item A function's local scope is determined by its definition, not by whichever function happened to call it.
    \item Calling \texttt{adjust} from \texttt{compute} does not make \texttt{compute}'s local names automatically local to \texttt{adjust}.
    \item Data should cross the call boundary explicitly through arguments and returned values.
  \end{itemize}

  \begin{block}{Preferred data flow}
    \centering
    caller value $\rightarrow$ argument $\rightarrow$ parameter $\rightarrow$ computation $\rightarrow$ return value
  \end{block}

  \begin{keyidea}
    Call-stack position describes control flow; lexical scope determines name visibility.
  \end{keyidea}
\end{frame}

\begin{frame}[shrink=4]{Worked Scope Trace: Three Different Bindings Named \texttt{x}}
  \begin{block}{Program}
    \texttt{def adjust(x: int) -> int:}\\
    \hspace*{0.20in}\texttt{y = x + 2}\\
    \hspace*{0.20in}\texttt{return y * 3}\\[0.05in]
    \texttt{def compute(value: int) -> int:}\\
    \hspace*{0.20in}\texttt{x = value + 1}\\
    \hspace*{0.20in}\texttt{answer = adjust(x)}\\
    \hspace*{0.20in}\texttt{return answer}\\[0.05in]
    \texttt{x = 10}\qquad \texttt{result = compute(4)}
  \end{block}

  \begin{center}
    \renewcommand{\arraystretch}{1.2}
    \begin{tabular}{l|l}
      \textbf{Environment} & \textbf{Relevant binding} \\\hline
      module level & \texttt{x -> 10} \\
      \texttt{compute} call & \texttt{x -> 5} \\
      \texttt{adjust} call & parameter \texttt{x -> 5}
    \end{tabular}
  \end{center}
\end{frame}

\begin{frame}{The Debugger Lets You Select Different Active Frames}
  \begin{columns}[T,onlytextwidth]
    \begin{column}{0.54\textwidth}
      \begin{itemize}
        \item When stopped inside a nested call, the call stack lists active frames.
        \item Selecting a frame changes which local bindings the debugger displays.
        \item The program has not changed; you are inspecting a different active environment.
      \end{itemize}

      \begin{activity}
        While stopped inside \texttt{adjust}, inspect the \texttt{adjust}, \texttt{compute}, and module-level frames. Identify every binding named \texttt{x}.
      \end{activity}
    \end{column}

    \begin{column}{0.40\textwidth}
      % Intended visual: IDE call-stack panel with three frames and separate variable panes.
      \lectureimage[2.70in]{slide14.png}
    \end{column}
  \end{columns}
\end{frame}

\sectionframe{Part 3: Function Composition}

\begin{frame}{Composition Connects Returned Values to New Computations}
  \begin{cpdefinition}
    \textbf{Function composition} uses the value produced by one function call as input to another computation.
  \end{cpdefinition}

  \begin{block}{Explicit data flow}
    \texttt{step1 = add\_one(4)}\\
    \texttt{result = double(step1)}\\[0.06in]
    Equivalent nested form: \texttt{result = double(add\_one(4))}
  \end{block}

  \begin{block}{Data-flow interpretation}
    \centering
    \texttt{4} $\rightarrow$ \texttt{add\_one} $\rightarrow$ \texttt{5}
    $\rightarrow$ \texttt{double} $\rightarrow$ \texttt{10}
  \end{block}
\end{frame}

\begin{frame}{Nested Calls Encode the Same Data Flow More Compactly}
  \begin{columns}[T,onlytextwidth]
    \begin{column}{0.46\textwidth}
      \begin{block}{With an intermediate name}
        \texttt{step1 = add\_one(4)}\\
        \texttt{result = double(step1)}
      \end{block}
    \end{column}

    \begin{column}{0.46\textwidth}
      \begin{block}{Nested composition}
        \texttt{result = double(add\_one(4))}
      \end{block}
    \end{column}
  \end{columns}

  \begin{keyidea}
    The nested form removes an intermediate binding; it does not remove the intermediate value or the need to evaluate it.
  \end{keyidea}
\end{frame}

\begin{frame}{Trace Nested Composition from the Inside Out}
  \begin{block}{Expression}
    \centering
    \Large\texttt{double(add\_one(4))}
  \end{block}

  \begin{enumerate}
    \item Identify the outer callable \texttt{double}.
    \item Evaluate its argument expression \texttt{add\_one(4)}.
    \item The \texttt{add\_one} body executes and returns \texttt{5}.
    \item The outer call now has argument value \texttt{5}.
    \item The \texttt{double} body executes and returns \texttt{10}.
  \end{enumerate}

  \begin{block}{Expression reduction}
    \centering
    \texttt{double(add\_one(4))} $\rightarrow$ \texttt{double(5)} $\rightarrow$ \texttt{10}
  \end{block}
\end{frame}

\begin{frame}{Composition Is Order-Sensitive}
  \begin{columns}[T,onlytextwidth]
    \begin{column}{0.46\textwidth}
      \begin{block}{Double, then square}
        \texttt{square(double(3))}\\[0.06in]
        \texttt{double(3) -> 6}\\
        \texttt{square(6) -> 36}
      \end{block}
    \end{column}

    \begin{column}{0.46\textwidth}
      \begin{block}{Square, then double}
        \texttt{double(square(3))}\\[0.06in]
        \texttt{square(3) -> 9}\\
        \texttt{double(9) -> 18}
      \end{block}
    \end{column}
  \end{columns}

  \begin{keyidea}
    Function composition is generally not commutative: changing the order can change the result.
  \end{keyidea}
\end{frame}

\begin{frame}{A Call May Contain Multiple Function-Call Arguments}
  \begin{block}{Example}
    \centering
    \Large\texttt{result = add(square(3), square(4))}
  \end{block}

  \begin{block}{Values required before \texttt{add} can execute}
    \texttt{square(3) -> 9}\\
    \texttt{square(4) -> 16}\\[0.06in]
    then \texttt{add(9, 16) -> 25}
  \end{block}

  \begin{activity}
    How many distinct calls occur? During which call is the parameter name \texttt{n} bound to \texttt{3}, and during which is it bound to \texttt{4}?
  \end{activity}
\end{frame}

\begin{frame}{Python Evaluates Argument Expressions from Left to Right}
  \begin{block}{Expression}
    \centering
    \texttt{add(square(3), square(4))}
  \end{block}

  \begin{enumerate}
    \item Evaluate the first argument expression \texttt{square(3)} $\rightarrow$ \texttt{9}.
    \item Evaluate the second argument expression \texttt{square(4)} $\rightarrow$ \texttt{16}.
    \item Bind \texttt{x -> 9} and \texttt{y -> 16} for the call to \texttt{add}.
    \item Execute \texttt{add} and return \texttt{25}.
  \end{enumerate}

  \begin{alertblock}{Tracing discipline}
    Do not merge the two calls to \texttt{square}; each call has its own parameter binding and lifetime.
  \end{alertblock}
\end{frame}

\sectionframe{Part 4: Tracing Functions That Call Functions}

\begin{frame}{Parameter Bindings Must Be Kept Separate Across Calls}
  \begin{block}{Program}
    \texttt{def difference(a: int, b: int) -> int:}\\
    \hspace*{0.20in}\texttt{return a - b}\\[0.06in]
    \texttt{def distance\_between(left: int, right: int) -> int:}\\
    \hspace*{0.20in}\texttt{distance = difference(right, left)}\\
    \hspace*{0.20in}\texttt{return distance}
  \end{block}

  \begin{center}
    \renewcommand{\arraystretch}{1.25}
    \begin{tabular}{l|l}
      \textbf{Call} & \textbf{Parameter bindings} \\\hline
      \texttt{distance\_between(4, 10)} & \texttt{left -> 4, right -> 10} \\
      \texttt{difference(right, left)} & \texttt{a -> 10, b -> 4}
    \end{tabular}
  \end{center}
\end{frame}

\begin{frame}{Full Trace: Follow Values Across Two Function Boundaries}
  \begin{enumerate}
    \item \texttt{distance\_between(4, 10)} creates \texttt{left -> 4}, \texttt{right -> 10}.
    \item Evaluate arguments to \texttt{difference}: \texttt{right -> 10}, \texttt{left -> 4}.
    \item The nested call creates \texttt{a -> 10}, \texttt{b -> 4}.
    \item \texttt{difference} returns \texttt{6}.
    \item The suspended outer call binds \texttt{distance -> 6}.
    \item \texttt{distance\_between} returns \texttt{6} to its own caller.
  \end{enumerate}

  \begin{keyidea}
    Record both the current frame and its bindings while tracing. Returned values travel one caller boundary at a time; they do not skip directly to the outermost code.
  \end{keyidea}
\end{frame}

\begin{frame}{A Return Always Resumes the Immediate Caller}
  \begin{columns}[T,onlytextwidth]
    \begin{column}{0.50\textwidth}
      \begin{block}{Call chain}
        initial code\\
        $\downarrow$\\
        \texttt{distance\_between}\\
        $\downarrow$\\
        \texttt{difference}
      \end{block}
    \end{column}

    \begin{column}{0.44\textwidth}
      \begin{block}{Return chain}
        \texttt{difference} returns\\
        $\downarrow$\\
        \texttt{distance\_between} resumes\\
        $\downarrow$\\
        then it may return to initial code
      \end{block}
    \end{column}
  \end{columns}

  \begin{alertblock}{Common mistake}
    A nested function's returned value is first delivered to the expression that called it, not automatically to every outer caller.
  \end{alertblock}
\end{frame}

\begin{frame}{Common Failure Modes with Nested Calls}
  \begin{center}
    \renewcommand{\arraystretch}{1.20}
    \begin{tabular}{p{0.43\textwidth}|p{0.49\textwidth}}
      \textbf{Incorrect model} & \textbf{Correction} \\\hline
      ``The caller disappears while the callee runs.'' & The caller remains active and suspended on the stack. \\
      ``Same variable name means same variable.'' & Bindings belong to scopes and execution environments. \\
      ``A callee can see the caller's locals.'' & Call relationships do not create lexical scope relationships. \\
      ``Nested syntax executes from the outside inward.'' & Argument values must be produced before the outer function body can execute. \\
      ``Returning exits every active function.'' & A return completes only the current call.
    \end{tabular}
  \end{center}
\end{frame}

\sectionframe{Part 5: Library Functions and Abstraction Boundaries}

\begin{frame}{Library Functions Participate in the Same Data-Flow Model}
  \begin{block}{Example}
    \texttt{from math import sqrt}\\[0.08in]
    \texttt{value = 25.0}\\
    \texttt{root = sqrt(value)}
  \end{block}

  \begin{itemize}
    \item Evaluate the argument expression \texttt{value} to obtain \texttt{25.0}.
    \item Call \texttt{sqrt} with that value.
    \item Receive the returned value \texttt{5.0}.
    \item Bind \texttt{root -> 5.0}.
  \end{itemize}

  \begin{keyidea}
    The same argument--return-value model applies even when we do not inspect the function's implementation.
  \end{keyidea}
\end{frame}

\begin{frame}[shrink=4]{Capstone Example: User-Defined Functions Plus \texttt{sqrt}}
  \begin{block}{Program}
    \texttt{from math import sqrt}\\[0.05in]
    \texttt{def square(n: float) -> float:}\\
    \hspace*{0.20in}\texttt{return n * n}\\[0.05in]
    \texttt{def magnitude(x: float, y: float) -> float:}\\
    \hspace*{0.20in}\texttt{x\_squared = square(x)}\\
    \hspace*{0.20in}\texttt{y\_squared = square(y)}\\
    \hspace*{0.20in}\texttt{result = sqrt(x\_squared + y\_squared)}\\
    \hspace*{0.20in}\texttt{return result}\\[0.05in]
    \texttt{answer = magnitude(3, 4)}
  \end{block}

  \begin{activity}
    Predict every intermediate value and the call stack while the first call to \texttt{square} is executing.
  \end{activity}
\end{frame}

\begin{frame}{Trace the Data Before Tracing the Implementation}
  \begin{columns}[T,onlytextwidth]
    \begin{column}{0.52\textwidth}
      \begin{block}{Data flow}
        \texttt{3 -> square -> 9}\\
        \texttt{4 -> square -> 16}\\
        \texttt{9 + 16 -> 25}\\
        \texttt{25 -> sqrt -> 5.0}\\
        \texttt{5.0 -> return -> answer}
      \end{block}

      \begin{itemize}
        \item First establish what values must flow between calls.
        \item Then inspect frames when a detailed execution trace is required.
      \end{itemize}
    \end{column}

    \begin{column}{0.40\textwidth}
      % Intended visual: data-flow graph for magnitude(3,4) with square and sqrt nodes.
      \lectureimage[2.85in]{slide30.png}
    \end{column}
  \end{columns}
\end{frame}

\begin{frame}{Use an Abstraction Boundary for Trusted Library Functions}
  \begin{columns}[T,onlytextwidth]
    \begin{column}{0.46\textwidth}
      \begin{block}{User-defined function under study}
        Trace inside when you need to understand:
        \begin{itemize}
          \item parameter bindings;
          \item local assignments;
          \item nested calls;
          \item returned values.
        \end{itemize}
      \end{block}
    \end{column}

    \begin{column}{0.46\textwidth}
      \begin{block}{Trusted library function}
        Often reason from its documented interface:
        \begin{itemize}
          \item accepted arguments;
          \item documented behavior;
          \item returned value.
        \end{itemize}
      \end{block}
    \end{column}
  \end{columns}

  \begin{keyidea}
    Abstraction lets us choose the appropriate level of detail for the reasoning task.
  \end{keyidea}
\end{frame}

\begin{frame}{Debugger Strategy: Step Into Selectively}
  \begin{itemize}
    \item Use \textbf{Step Into} when the implementation of a user-defined call is part of the concept you are tracing.
    \item Use \textbf{Step Over} when a trusted call can be treated through its interface.
    \item Inspect the call stack whenever more than one user-defined call is active.
    \item Select different frames to compare their parameter and local-variable bindings.
  \end{itemize}

  \begin{block}{Lab 3 discipline}
    Predict the stack and bindings first; then use the debugger to test the prediction.
  \end{block}
\end{frame}

\sectionframe{Part 6: Decomposition and Program Design}

\begin{frame}{Decomposition and Composition Are Complementary Ideas}
  \begin{columns}[T,onlytextwidth]
    \begin{column}{0.46\textwidth}
      \begin{cpdefinition}
        \textbf{Decomposition} separates a larger computation into smaller functions with focused responsibilities.
      \end{cpdefinition}
    \end{column}

    \begin{column}{0.46\textwidth}
      \begin{cpdefinition}
        \textbf{Composition} reconnects those functions by passing produced values into later computations.
      \end{cpdefinition}
    \end{column}
  \end{columns}

  \begin{block}{Design pattern}
    \centering
    large problem $\rightarrow$ smaller transformations $\rightarrow$ connected data flow
  \end{block}
\end{frame}

\begin{frame}{Worked Design: Separate Discount and Tax Calculations}
  \begin{columns}[T,onlytextwidth]
    \begin{column}{0.47\textwidth}
      \begin{block}{Smaller transformations}
        \texttt{def apply\_discount(price, rate):}\\
        \hspace*{0.18in}\texttt{return price - price * rate}\\[0.07in]
        \texttt{def apply\_tax(price, rate):}\\
        \hspace*{0.18in}\texttt{return price + price * rate}
      \end{block}
    \end{column}

    \begin{column}{0.47\textwidth}
      \begin{block}{Composition}
        \texttt{def final\_price(price, discount, tax):}\\
        \hspace*{0.18in}\texttt{discounted = apply\_discount(price, discount)}\\
        \hspace*{0.18in}\texttt{return apply\_tax(discounted, tax)}
      \end{block}
    \end{column}
  \end{columns}

  \begin{keyidea}
    Each function names one transformation; the outer function expresses how the transformations are connected.
  \end{keyidea}
\end{frame}

\begin{frame}{Good Interfaces Make Composition Easier to Reason About}
  \begin{itemize}
    \item Parameters should represent the data a computation genuinely requires.
    \item Returned values should expose the result needed by the caller.
    \item Avoid relying on unrelated external state when a value can be passed explicitly.
    \item Prefer functions whose names communicate the transformation they perform.
  \end{itemize}

  \begin{block}{Mental contract}
    \centering
    \textbf{required input data} $\rightarrow$ \textbf{named computation} $\rightarrow$ \textbf{returned result}
  \end{block}

  \begin{alertblock}{Why this matters for tracing}
    Explicit inputs and outputs make the path of data visible across function boundaries.
  \end{alertblock}
\end{frame}

\begin{frame}[shrink=5]{Practice: Trace a Composed Program}
  \begin{block}{Program}
    \texttt{def triple(n: int) -> int:}\\
    \hspace*{0.20in}\texttt{return n * 3}\\[0.05in]
    \texttt{def difference(a: int, b: int) -> int:}\\
    \hspace*{0.20in}\texttt{return a - b}\\[0.05in]
    \texttt{def compute(x: int, y: int) -> int:}\\
    \hspace*{0.20in}\texttt{left = triple(x)}\\
    \hspace*{0.20in}\texttt{right = triple(y)}\\
    \hspace*{0.20in}\texttt{return difference(left, right)}\\[0.05in]
    \texttt{answer = compute(5, 2)}
  \end{block}

  \begin{activity}
    Predict: (1) every parameter binding, (2) every returned value, (3) the maximum call-stack depth, and (4) the final binding of \texttt{answer}.
  \end{activity}
\end{frame}

\begin{frame}[shrink=12]{Takeaways}
  \begin{itemize}
    \item A function may begin another call before it has finished; unfinished callers remain active.
    \item The call stack records active calls, while each call has its own runtime execution environment.
    \item Scope determines which binding a name refers to; call-stack position does not create new lexical visibility.
    \item Returned values let functions be composed so that one computation supplies data to another.
    \item Nested argument expressions must produce values before the receiving function body can execute.
    \item Returning completes only the current call and resumes its immediate caller.
    \item Library functions can often be reasoned about through their documented input--output interface.
    \item Decomposition creates focused transformations; composition connects them into a larger program.
  \end{itemize}

  \begin{keyidea}
    Mental model: \textbf{values flow through calls; active calls keep separate bindings; return values reconnect the suspended computation}.
  \end{keyidea}
\end{frame}

% Reference notes for instructor/source maintenance:
% Python Language Reference, Expressions: calls, return values, and left-to-right evaluation order.
% Python Language Reference, Execution model: name binding, local/global scope, and name resolution.

\end{document}
